\section{{\fontsize{12pt}{14pt}\selectfont INTRODUÇÃO}}
    
\vspace{-0.3cm}
\novoBlocoTexto {
    Ao estudar qualquer conteúdo, é inevitável que os estudantes encontrem ambiguidades ou dificuldades e precisem de auxílio \parencite{ryan2001avoiding}. Nem toda dúvida, porém, é levada ao professor. Dillon (1988, apud \cite{ryan2001avoiding}) chama atenção para a quantidade de perguntas que permanecem silenciadas, destacando que muitas vezes se prefere evitar a verbalização e aparentar compreensão. 
}
\novoBlocoTexto {
    No ensino superior, muitos estudantes deixam de buscar a ajuda de que precisam por considerá-la constrangedora ou uma admissão de fracasso \parencite{karabenick2011understanding}. Entre os principais motivos para evitar perguntas, destacaram-se a incerteza sobre a qualidade da questão formulada e o receio do julgamento dos colegas. 
}
\novoBlocoTexto {
    Os próprios professores percebem esse silêncio. Em entrevistas com professoras do Ensino Fundamental, \textcite{mertins2022valorizacao} registraram que a vergonha, a timidez e o medo de se expor diante da turma inibem as perguntas, e uma das professoras resumiu: "Se tu só pergunta, alguém tem alguma dúvida? [...] Ninguém, né?". Embora as pesquisas apresentadas tenham sido realizadas em diferentes contextos educacionais, seus resultados apontam para a presença de fatores sociais e emocionais que podem dificultar a manifestação de dúvidas, independentemente das características específicas dos conteúdos estudados.
}
\novoBlocoTexto {
    Para o docente, essa lacuna tem consequências diretas. As professoras entrevistadas por \textcite{mertins2022valorizacao} relatam que as perguntas dos estudantes permitem identificar o que eles sabem e o que ainda não compreenderam, contribuindo para o planejamento das aulas. Uma das participantes afirma: "A partir do momento que eles trazem as perguntas, elaboro a aula" \space \parencite[p.~5]{mertins2022valorizacao}. Quando essas dúvidas não são manifestadas, o professor pode identificar tardiamente as diferenças entre o que foi ensinado e o que foi efetivamente aprendido, o que reduz as possibilidades de intervenção pedagógica \parencite{angelo1993classroom}. 
}
\novoBlocoTexto{
    Além disso, nem toda dúvida que deixa de chegar ao professor permanece sem resposta. Muitos estudantes recorrem inicialmente aos colegas, aos livros ou à internet e procuram o docente apenas quando a dificuldade persiste \parencite{nguyenviet2024exploratory}. Com o avanço das ferramentas de inteligência artificial generativa, surgiram novas possibilidades para o esclarecimento dessas dúvidas. Em relato de experiência com turmas iniciantes de programação que utilizaram o ChatGPT e o DeepSeek, uma das professoras observou que a maioria dos estudantes não buscou ajuda, pois as ferramentas esclareciam suas dúvidas de maneira rápida e eficaz, e manifestou preocupação com a possível dependência desses recursos \parencite{maia2025utilizacao}.
}
\novoBlocoTexto {
    Em pesquisa com estudantes de graduação do Reino Unido, a proporção dos que utilizavam inteligência artificial generativa para obter explicações sobre conceitos passou de 36\%, em 2024, para 58\%, em 2025, mantendo-se como a finalidade de uso mais frequente entre os participantes \parencite{freeman2025student}. No Brasil, 70\% dos estudantes do ensino médio já utilizavam inteligência artificial generativa em pesquisas escolares em 2024 \parencite{nicbr2025tic}. Esses dados não apontam, por si só, uma redução na procura pelos professores, mas sugerem que uma parcela crescente das dúvidas passa a ser tratada em canais aos quais o docente não tem acesso. 
}
\novoBlocoTexto{
    Dessa forma, parte das dificuldades dos estudantes pode permanecer desconhecida pelo professor, seja pela resistência em verbalizar dúvidas, seja pelo uso de fontes alternativas para auxílio. Ainda que esses recursos contribuam para a aprendizagem individual, a falta de informações sobre as dificuldades encontradas limita a compreensão do professor acerca das necessidades de sua turma. 
}
\novoBlocoTexto{
    A partir dessa problemática, estabelece-se a seguinte pergunta de pesquisa: Como proporcionar ao docente acesso organizado às dúvidas dos estudantes, com base em seu próprio material didático, de modo a subsidiar o aprimoramento de suas práticas de ensino?
}
\novoBlocoTexto{
    Para responder a essa questão, o presente trabalho tem como objetivo geral desenvolver e avaliar um protótipo web de um assistente educacional baseado em LLM (sigla em inglês para \textit{Large Language Model}, ou Modelos de Linguagem de Grande Escala) e RAG (sigla em inglês para \textit{Retrieval-Augmented Generation}, ou Geração Aumentada por Recuperação) que classifica e organiza as dúvidas de estudantes com base no plano de ensino do próprio docente, disponibilizando essas informações em um painel para acompanhamento dos conteúdos questionados pelos estudantes.
}